\scanpage{85}
% archive foliation: 84
% lines 1--4
se déduit de l'hom évident
\[
\mu_n(\overline Q)\xrightarrow{\rho_\nu}\widetilde A_\nu
\]
provenant de la permutation circulaire des
gén. de $\widehat L_\nu$ -- l'opération de
% lines 5--8
$\mathbb Z/_{2}$ dans $\widetilde{\unsure{\mathcal T|\Gamma}}_\nu$ \unsure{provient|provenant} de $=$ de
l'opération
\[
\mathbb Z/_{2}\text{---}\widetilde A_\nu
\]
\deleted{\unsure{par}}
% lines 9--11
\margin{\uline{vérifier} ?}
\[
\left\{\begin{aligned}
(\lambda_0,\lambda_1,\ldots\lambda_\nu)\xrightarrow{t_\nu}&\deleted{(\unsure{\lambda_0},\lambda_\nu,\lambda_{\nu-1},\ldots\lambda_1)}\\
&\inserted{(\lambda_0,q_\nu\lambda_\nu q_\nu^{-1},q_{\nu-1}\lambda_{\nu-1}q_{\nu-1}^{-1},\ldots q_2\lambda_2q_2^{-1},\lambda_1)}\\
\text{où }&\inserted{\unsure{q_i|g_i}=\lambda_{\unsure{0}}\cdots\lambda_{i-1}}
\end{aligned}\right.
\]
qui renverse l'ordre circulaire.
% lines 12--14
Ceci nous dit donc que l'hom
\diagram{\begin{tikzcd}[ampersand replacement=\&]
\Gamma\arrow[r,"i_\nu",dash]\&\widetilde{\unsure{\mathcal T|\Gamma}}_\nu\arrow[r,"\varepsilon_\nu",dotted,dash]\&\unsure{\mathfrak S_n|\mathcal S_n}\supset(\mathbb Z/_{n})^{*}
\end{tikzcd}}
\deleted{$\unsure{\mu_n(\overline Q)}$}
% lines 15--17
satisfait, en plus de
\[
\left[\varepsilon_\nu i_\nu=\chi_n\right.\quad\unsure{on|ou}
\]
la relation
\[
\left\{\begin{aligned}
&i_\nu(\Gamma)\text{ commute à }t_\nu\\
&i_\nu(\gamma)\rho_\nu(i)\unsure{i_\nu|t_\nu}(\gamma)^{-1}=\rho_\nu(\chi_n(\gamma)i)
\end{aligned}\right.
\]
